% ============================================================================
%  Discussion / outlook paragraph: where a shear-dependent outlet closure
%  earns its keep.
%
%  Two versions.  The short one is the paragraph; the long one splits it in
%  three and adds the constitutive-uncertainty argument, which is the part a
%  referee is most likely to press on.  Use one or the other, not both.
% ============================================================================

% ---------------------------------------------------------------- short ----
The present results delimit rather than diminish the role of a shear-dependent
outlet closure. Throughout the perfused coronary bed the reduced shear rate
satisfies $\lambda\dot\gamma \gg 1$ --- the identity
$\dot\gamma_0 = r\,\Delta p/(2\mu_N L)$ places every compartment between
$10^2$ and $10^3\,\mathrm{s}^{-1}$ --- so the constitutive law is evaluated on
its high-shear plateau for essentially the whole cycle and the apparent
viscosity departs from $\mu_\infty$ by only a few per cent. A stenosis does not
relieve this constraint but tightens it: an $80\,\%$ reduction in area raises
the throat shear rate by two orders of magnitude and drives the fluid closer
still to Newtonian behaviour. The pathological signal in hyperglycaemia and in
hyperviscosity therefore reaches the flow through the \emph{level} of the
plateau --- haematocrit and plasma composition --- and not through the
shear-thinning branch of the law. Closures of the present type should instead
become consequential wherever the shear rate falls to
$O(1\text{--}10\,\mathrm{s}^{-1})$: in ectatic and aneurysmal coronary
segments, where $\dot\gamma \sim Q/r^3$ makes a $1.5$--$3$-fold segmental
dilatation cost one to one and a half orders of magnitude of shear, and, with
a wider margin, in the fibrillating atrial appendage, the venous valve sinus
and the aneurysmal aortic sac.

% ----------------------------------------------------------------- long ----
% The present results delimit rather than diminish the role of a shear-dependent
% outlet closure. Throughout the perfused coronary bed the reduced shear rate
% satisfies $\lambda\dot\gamma \gg 1$ --- the identity
% $\dot\gamma_0 = r\,\Delta p/(2\mu_N L)$ places every compartment between
% $10^2$ and $10^3\,\mathrm{s}^{-1}$ --- so the constitutive law is evaluated on
% its high-shear plateau for essentially the whole cycle and the apparent
% viscosity departs from $\mu_\infty$ by only a few per cent. A stenosis does not
% relieve this constraint but tightens it: an $80\,\%$ reduction in area raises
% the throat shear rate by two orders of magnitude and drives the fluid closer
% still to Newtonian behaviour. The pathological signal in hyperglycaemia and in
% hyperviscosity therefore reaches the flow through the \emph{level} of the
% plateau --- haematocrit and plasma composition --- and not through the
% shear-thinning branch of the law.
%
% Closures of the present type should become consequential wherever the shear
% rate falls to $O(1\text{--}10\,\mathrm{s}^{-1})$. Coronary artery ectasia,
% reported in $1$--$5\,\%$ of diagnostic angiograms, is the most direct such
% setting within the coronary circulation: since $\dot\gamma \sim Q/r^3$, a
% $1.5$--$3$-fold segmental dilatation costs one to one and a half orders of
% magnitude of shear, and the angiographic hallmarks of the condition ---
% delayed antegrade filling, segmental backflow, contrast stasis --- are the
% signature of a recirculating flow of long residence time. That coronary flow
% reserve is depressed in ectatic segments in the absence of any obstruction
% indicates an impairment carried by the flow field rather than by a luminal
% cross-section, which is precisely what a lumped shear-dependent impedance is
% built to represent. The same argument extends, with a wider margin, to the
% fibrillating atrial appendage, the venous valve sinus and the aneurysmal
% aortic sac.
%
% Two cautions attach to that extrapolation. The first is quantitative: the
% excess apparent viscosity predicted at $\dot\gamma \approx 10^2\,\mathrm{s}^{-1}$
% differs by more than a factor of two between published Carreau--Yasuda
% parameterizations of human blood, so at conduit shear rates the constitutive
% uncertainty exceeds the modelled effect, whereas below $10\,\mathrm{s}^{-1}$
% the same parameterizations agree that the excess is of order unity. Low-shear
% applications are therefore not merely larger in effect but better posed. The
% second is structural: as the residence time grows towards the seconds-long
% timescale of erythrocyte aggregation, thixotropy and a finite yield stress
% cease to be negligible, and the equilibrium relation $\mu(\dot\gamma)$ assumed
% here must give way to an evolution equation for the aggregate structure. The
% regime in which the present closure is most valuable is thus also the regime
% in which it is closest to its own limit of validity.
